\section{能效问题栈：FLOPS 之外还有三堵墙}

从应用扩散回到底层，首先遇到的是能效。讲者列出 semiconductor、memory、interconnect 与 cooling 四个方向，实际想表达的是：AI system 的吞吐受最慢资源限制，峰值 FLOPS 只是一个上界。本节用 roofline 式思维把这些约束统一起来，避免把所有瓶颈都叫“缺 GPU”。

\subsection{Von Neumann bottleneck 与 data movement}

\term{von Neumann bottleneck}（冯·诺依曼瓶颈）指处理器执行单元与存储分离，数据和指令必须在 memory hierarchy 与 compute 之间移动。许多 AI kernel 的能量和时间并不消耗在乘加本身，而消耗在读取 weights、KV cache、activations 和跨设备交换。对一个 workload，可写出简化下界
\[
T_{step}\geq \max\left(\frac{F}{C},\frac{B_m}{BW_m},\frac{B_n}{BW_n}\right),
\]
其中 $F$ 是所需运算量，$C$ 是有效计算吞吐；$B_m$ 与 $BW_m$ 分别是 memory traffic 和 memory bandwidth；$B_n$ 与 $BW_n$ 是 network traffic 与 interconnect bandwidth。公式说明：只提高 $C$，另外两项仍可能主导时间。

\lecturefigure{02-efficiency-stack.png}{AI 能效需要同时检查 semiconductor、memory、interconnect 和 facility。}{本地字幕 06:20--12:00；概念重绘。}

\begin{knowledgebox}{读图：四层不是四个独立项目}
新 semiconductor 可能提高 rack density，进而改变 cooling 和供电；更大 \term{High Bandwidth Memory}（HBM，高带宽堆叠 DRAM）可能减少跨卡通信，却提高封装成本；optical interconnect 可能降低长距离传输损失，却引入新器件和软件栈。正确问题是端到端 workload 获得多少 SLO-compliant tokens per watt，而不是某个器件的峰值指标是否更漂亮。
\end{knowledgebox}

\subsection{Compound semiconductor 解决什么，不解决什么}

\term{compound semiconductor}（化合物半导体）由两种或多种元素组成，例如 GaN（gallium nitride，氮化镓）或 SiC（silicon carbide，碳化硅）。它们在高电压、高频和功率转换中可提供更好的击穿场强或效率，因此适合电源、射频和光电器件；但“从 silicon 全面迁移”不是一个单一开关。逻辑芯片、memory、power electronics 与 photonics 对材料、制造良率和成本的要求不同。

\begin{warningbox}{课堂中的器件比喻不能当产品路线图}
字幕把 switching voltage、leakage 和 compound materials 混在一个快速回答里。讲义保留“器件层仍有能效空间”的判断，但不把某个材料直接宣称为通用 GPU 替代品。技术可行性还取决于 process integration、yield、packaging、toolchain 与规模经济。
\end{warningbox}

\subsection{本章小结}

AI 基础设施的能效是端到端系统属性。算术、memory、network 和 facility 任何一层不足，都可能让昂贵 accelerator 等待；研究与采购都应以真实 workload 的 useful output 为共同指标。

